\section{Introduction}
\label{sec:introduction}

Millimeter-wave links buy their bandwidth from carrier frequencies at which
oscillator phase noise is severe. A 60~GHz synthesizer built to a commercial
power budget has a 3~dB corner linewidth in the hundreds of kilohertz, and the
phase it contributes wanders by several degrees over a single orthogonal
frequency division multiplexing (OFDM) symbol at the symbol durations these
systems use~\citep{nakagome2019oscillator}.

In a fully digital array this would be a manageable problem, because
independent chains carry independent oscillators and the impairment partially
averages across antennas. Hybrid arrays do not have that property. One local
oscillator drives the entire analog combining network, so the phase-noise
process multiplies the combined signal after beamforming and appears identically
on every spatial stream~\citep{oyelaran2020hybrid}. Array gain does nothing to
it.

The conventional treatment splits phase noise into a common phase error, a
single rotation shared by all subcarriers within a symbol, and intercarrier
interference (ICI), which is discarded as
noise~\citep{bertrand2018commonphase}. That split is sound when the linewidth
times the symbol duration is small. At 400~MHz of bandwidth and a 200~kHz
linewidth it is not: the ICI carries more energy than the residual thermal
noise at operating signal-to-noise ratios, and, critically, it is not white.
Treating structured interference as white noise is the specific failure this
paper addresses.

\paragraph*{Contributions}
We make three.

First, we derive the exact covariance of the ICI term for a free-running
oscillator with a Lorentzian spectrum, without the small-phase approximation
that the literature usually applies at this point
(Section~\ref{sec:system-model}). The covariance is Toeplitz and effectively
rank deficient, and we characterise its numerical rank as a function of
linewidth and symbol duration.

Second, we build \est, a variational estimator that treats the phase trajectory
as a latent Gauss-Markov process and alternates between a channel update in
closed form and a phase update by Kalman smoothing
(Section~\ref{sec:estimator}). The complexity is $O(N \log N + N R^2)$ per
iteration for $N$ subcarriers and numerical rank $R$, against $O(N^3)$ for a
direct whitened least squares solution (Section~\ref{sec:complexity}).

Third, we evaluate against three baselines over a linewidth sweep spanning two
orders of magnitude, and report where the method stops helping
(Section~\ref{sec:results}).
